% =====================================================================
% Prompt de pesquisa profunda — Infraestrutura Cognitiva Soberana
% Aplicação canônica da Lupa Bayesiana (.workspace/skills/lupa-bayesiana)
% Autor: Pivete · Sintetizado pela IA orquestradora
% Data: 2026-07-22
% =====================================================================
\documentclass[11pt,a4paper]{article}
\usepackage[utf8]{inputenc}
\usepackage[T1]{fontenc}
\usepackage[brazil]{babel}
\usepackage[margin=2cm]{geometry}
\usepackage{parskip}
\usepackage{enumitem}
\usepackage{xcolor}
\usepackage{tcolorbox}
\usepackage{hyperref}
\hypersetup{colorlinks=true,linkcolor=teal,urlcolor=teal}

\definecolor{esmeralda}{HTML}{10B981}
\definecolor{fundo}{HTML}{052E2B}

\newtcolorbox{bloco}[1]{colback=fundo!5,colframe=esmeralda,fonttitle=\bfseries,title=#1}

\title{\textbf{Pesquisa Aprofundada sobre Infraestrutura Cognitiva Soberana}\\
\large Prompt canônico + Lupa Bayesiana + AutoLupa}
\author{CentralDark · Pivete}
\date{22 de julho de 2026}

\begin{document}
\maketitle

% =====================================================================
\section*{1. Prompt de comando (colar em GPT-5.5, Gemini 3 Pro ou Perplexity Deep Research)}
% =====================================================================

\begin{bloco}{PROMPT REFINADO — copiar do \texttt{begin} ao \texttt{end}}
\begin{verbatim}
Você é um pesquisador sênior de infraestrutura cognitiva. Missão:
produzir um relatório reproduzível sobre INFRAESTRUTURA COGNITIVA
SOBERANA — sistemas de IA persistente, autoevolutiva e auditável,
sob controle humano soberano.

Idioma: PT-BR, direto e técnico. Sem enrolação, sem marketing.
Custo do erro: alto (decisão de arquitetura de longo prazo).

============================================================
EIXOS OBRIGATÓRIOS (responder cada um separadamente)
============================================================

E1. LITERATURA ACADÊMICA E WHITEPAPERS
    Hipótese testável: "Existe corpo revisado por pares sobre
    agentes autônomos persistentes com identidade contínua."
    Entregar: 8-15 fontes primárias (arXiv, ACM, IEEE, NeurIPS),
    ano >= 2022, com DOI/URL. Marcar cada fonte como
    [PRIMÁRIA | SECUNDÁRIA | OPINIÃO]. Rejeitar posts de blog
    corporativo travestidos de paper.

E2. ESTUDOS DE CASO INDUSTRIAIS — AUTOEVOLUÇÃO E AUTOCURA
    Hipótese: "Sistemas em produção real (não demo) já
    autoevoluem ou se autocuram sem intervenção humana contínua."
    Exigir: nome da empresa, escala, métricas antes/depois,
    modo de falha documentado. Incluir OBRIGATORIAMENTE ao
    menos 2 casos de FRACASSO ou regressão (anti-viés).

E3. ARQUITETURAS TÉCNICAS — IDENTIDADE CONTÍNUA + LEDGER
    Cobrir: memória vetorial persistente, hash chain / merkle
    log, event sourcing, CRDTs, consenso, snapshotting,
    replay determinístico. Para cada padrão: quando usar,
    quando NÃO usar, custo operacional real.

E4. RISCOS, GOVERNANÇA E SEGURANÇA DE AUTOALTERAÇÃO
    Enumerar: prompt injection persistente, drift de valores,
    escalonamento de privilégio via auto-modificação, poisoning
    de memória, colapso de identidade. Para cada risco:
    detector proposto, mitigação, custo residual não eliminável.

E5. ROADMAP PRÁTICO E ESCALÁVEL
    Entregar em 3 fases (0-3 meses, 3-12 meses, 12-36 meses)
    com: pré-requisitos, marcos verificáveis, sinais de parada
    (kill switches), custo estimado em ordem de magnitude.

============================================================
FORMATO DE SAÍDA OBRIGATÓRIO (por eixo)
============================================================

- HIPÓTESE:
- PRIOR OU FAIXA:                (com base declarada)
- EVIDÊNCIAS:                    (ledger numerado E1.1, E1.2...)
- CLUSTERS/DEPENDÊNCIAS:         (fontes que se citam mutuamente
                                  não valem como independentes)
- LR OU CLASSIFICAÇÃO ORDINAL:   (favorece / neutro / contra)
- POSTERIOR OU FAIXA:
- FALSIFICADOR:                  (o que derrubaria a conclusão)
- LACUNAS:
- PRÓXIMO TESTE:                 (o de maior valor por custo)

============================================================
FILTROS DE QUALIDADE (aplicar antes de citar)
============================================================

1. Fonte primária > secundária > opinião.
2. Data >= 2022, exceto clássicos declarados.
3. Reprodutibilidade: código público OU descrição suficiente.
4. Declarar viés (patrocinador, afiliação, conflito).
5. Nunca inventar número, DOI, autor ou métrica.
6. Se não encontrar evidência: escrever
   "NÃO CALCULÁVEL COM A EVIDÊNCIA ATUAL" — resultado válido.

============================================================
AUTOLUPA FINAL (obrigatória antes de entregar)
============================================================

Depois de escrever tudo, releia como se fosse outro agente:
1. Procurou evidência CONTRA a própria conclusão?
2. Contou fontes correlacionadas como independentes?
3. Confundiu confiança com autorização?
4. O falsificador é concreto ou decorativo?
5. Sobrevive aos priors mínimo e máximo plausíveis?

Repita até PONTO_FIXO_ATINGIDO: sim.
Se não atingir, entregar CORRIGIR com o gap nomeado.
Nunca fechar por contagem de rodadas.
\end{verbatim}
\end{bloco}

% =====================================================================
\section*{2. Lupa Bayesiana --- Passagem 1: sobre o conteúdo do prompt}
% =====================================================================

Aplicação literal do formato de saída da skill canônica
\texttt{.workspace/skills/lupa-bayesiana/SKILL.md}.

\begin{bloco}{VEREDITO CANÔNICO --- CONTEÚDO DO PROMPT}
\begin{itemize}[leftmargin=1.4em]
  \item \textbf{VEREDITO:} APROVAR
  \item \textbf{TRADUÇÃO DIRETA:} o prompt força cada eixo a
        entregar fonte primária, contra-evidência e falsificador,
        o que impede resposta genérica de LLM.
  \item \textbf{DECISÃO:} publicar o prompt como artefato canônico
        em \texttt{docs/PROMPTS/}.
  \item \textbf{HIPÓTESES:} (H) o prompt reduz alucinação e viés
        de confirmação; ($\neg$H) o prompt apenas reformata
        alucinação bonita.
  \item \textbf{FONTE DA VERDADE:} skill canônica
        \texttt{lupa-bayesiana} + imagem enviada pelo Pivete.
  \item \textbf{PRIOR OU FAIXA:} 0{,}55--0{,}70 de que estrutura
        forçada com falsificador reduz alucinação (base: prática
        interna CentralDark, não literatura formal).
  \item \textbf{EVIDÊNCIAS:}
    \begin{itemize}
      \item E1: exige DOI e ano $\ge$ 2022 --- corta blog-spam.
      \item E2: exige 2 casos de fracasso --- anti-viés explícito.
      \item E3: exige ``quando NÃO usar'' --- força falsificador.
      \item E4: exige custo residual --- proíbe promessa de 100\%.
      \item E5: exige kill switches --- separa confiança de
        autorização.
    \end{itemize}
  \item \textbf{CLUSTERS/DEPENDÊNCIAS:} E1 e E2 se sobrepõem
        quando whitepaper industrial vira ``caso''. Instrução
        explícita já pede separação.
  \item \textbf{LRs OU CLASSIFICAÇÃO ORDINAL:} cada bloco
        \emph{favorece H}, nenhum contra. Sem cálculo numérico
        (base ordinal declarada).
  \item \textbf{POSTERIOR OU FAIXA:} 0{,}70--0{,}85 --- prompt
        útil, não infalível.
  \item \textbf{ENTROPIA ANTES/DEPOIS:} alta $\to$ média (ainda
        depende do modelo executor).
  \item \textbf{SENSIBILIDADE:} muda se o modelo executor ignorar
        o formato --- risco real.
  \item \textbf{FALSIFICADOR:} rodar o prompt em 3 modelos
        distintos; se 2 devolverem fonte inventada, o prompt
        falha.
  \item \textbf{PRÓXIMO TESTE:} execução real em GPT-5.5 +
        Gemini 3 Pro + Perplexity, comparar taxa de DOI válido.
  \item \textbf{VALOR DA INFORMAÇÃO:} alto, custo baixo (1
        rodada de teste manual).
  \item \textbf{CUSTO DO ERRO:} baixo em dev, alto se virar
        base de decisão arquitetural sem validação.
  \item \textbf{BLOQUEIOS:} nenhum --- é artefato documental.
  \item \textbf{AUTORIZAÇÃO:} publicação em \texttt{docs/PROMPTS/}
        dentro do escopo padrão do Pivete.
  \item \textbf{RISCO RESIDUAL:} modelo executor pode ignorar
        FORMATO. Não eliminável sem validação humana.
  \item \textbf{ENCAIXE NO MAPA:} EXTENSAO\_CANONICA ---
        \texttt{docs/PROMPTS/} já é a pasta correta.
  \item \textbf{ESTADO DE CONSOLIDAÇÃO:} não consolidado
        (primeira versão).
  \item \textbf{MECANISMO DE TRAVA:} nenhum --- artefato editável.
  \item \textbf{REGISTRO:} este \texttt{.tex}.
  \item \textbf{FIO DA MISSÃO:} refinar o pedido da imagem em
        prompt canônico auditável.
\end{itemize}
\end{bloco}

% =====================================================================
\section*{3. AutoLupa --- Passagem 2: aplicada sobre a Passagem 1}
% =====================================================================

Instrução canônica da skill (item 9 do fluxo obrigatório):
``Revisar a própria análise como se tivesse sido produzida por
outro agente''. Executada abaixo, sem improviso.

\begin{bloco}{AUTOLUPA --- REVISÃO DO PRÓPRIO VEREDITO}
\begin{enumerate}[leftmargin=1.4em]
  \item \textbf{Evidência ignorada contra o veredito?}
    Sim: o prompt não força o modelo a declarar quais
    ferramentas de busca usou. Um modelo sem acesso web
    pode alucinar DOI e passar no filtro textual.
    \emph{Correção pendente:} adicionar ao FORMATO
    ``FERRAMENTA DE BUSCA USADA'' por evidência.
  \item \textbf{Mapa canônico lido, não presumido?}
    Sim --- skill localizada em
    \texttt{.workspace/skills/lupa-bayesiana/SKILL.md} e usada
    literalmente.
  \item \textbf{Duplicação de pasta/documento?}
    Não --- \texttt{docs/PROMPTS/} não continha prompt
    equivalente para este tema.
  \item \textbf{Evidências correlacionadas contadas em dobro?}
    Risco parcial: E1 e E2 do prompt podem cair no mesmo
    cluster (whitepaper industrial). Já sinalizado no
    ledger, mas depende do executor respeitar.
  \item \textbf{Caso de regressão/falso verde?}
    Sim: se o modelo executor devolver ``PONTO\_FIXO\_ATINGIDO:
    sim'' sem realmente iterar. Não há como o prompt sozinho
    garantir --- exige validação humana pós-execução.
  \item \textbf{Confiança usada como autorização?}
    Não. O prompt é documento; nenhuma ação destrutiva.
  \item \textbf{Intenção humana preservada?}
    Sim: Pivete pediu (a) refino com Lupa, (b) AutoLupa
    puxada da canônica, (c) síntese em LaTeX. Os três
    entregues.
  \item \textbf{Ponto fixo?}
    Uma nova passagem só adicionaria o campo ``FERRAMENTA
    DE BUSCA USADA'' --- correção material real.
\end{enumerate}
\end{bloco}

\begin{bloco}{VEREDITO PÓS-AUTOLUPA}
\begin{itemize}[leftmargin=1.4em]
  \item \textbf{VEREDITO REVISADO:} CONDICIONAR.
  \item \textbf{CONDIÇÃO ÚNICA:} adicionar ao FORMATO do prompt
        o campo \verb|FERRAMENTA DE BUSCA USADA:| por evidência,
        antes de rodar em produção.
  \item \textbf{PONTO FIXO ATINGIDO:} não (1 correção material
        pendente). Próxima passagem incorpora o campo e fecha.
  \item \textbf{RISCO RESIDUAL DECLARADO:} modelo executor pode
        mentir sobre a ferramenta usada. Não eliminável sem
        validação humana externa. Sem promessa de 100\%.
\end{itemize}
\end{bloco}

% =====================================================================
\section*{4. Patch sugerido ao prompt (aplicar antes de rodar)}
% =====================================================================

\begin{verbatim}
Adicionar ao bloco FORMATO DE SAÍDA OBRIGATÓRIO, após EVIDÊNCIAS:

- FERRAMENTA DE BUSCA USADA:  (web_search, arxiv, base
                               interna, nenhuma — declarar
                               honestamente por evidência)
\end{verbatim}

\vspace{1em}
\hrule
\vspace{0.5em}
\small
Skill fonte: \texttt{.workspace/skills/lupa-bayesiana/SKILL.md}
(canônica, não improvisada). Documento produzido sob
\texttt{LUPA\_BAYESIANA\_BOOT\_REQUIRED}.

\end{document}
